\section{Evaluation}
\label{sec:evaluation}

\subsection{Experimental Setup}

All 2020 GeoSection results use 2020 Census data, TIGER tract geometries,
and 2020 presidential election results (precinct-to-tract interpolation via
VEST \citep{fekrazad2021}).
District counts follow the 2020 apportionment.
All runs use 50 seeds per ratio.
Population balance tolerance is $0.5\,\%$, matching the congressional standard
from \emph{Karcher v.\ Daggett} (1983).

The reference implementation is the \texttt{redist} Rust binary, partition
mode \texttt{geosection}, using METIS 5.1 via the \texttt{metis-rs} crate.
All 44 multi-district states converge within the $0.5\,\%$ balance tolerance.

For 2000 census year data, the 50-state sweep has completed 47 states;
3 large states (CA, TX, FL) are still processing.
\textbf{[TBD: 2010 sweep pending]}: the 2010 census sweep is in progress and
will provide the third data point for three-census CSS computation.

\subsection{2020 Baseline: Full 50-State Natural Ratio Table}
\label{sec:2020-baseline}

Table~\ref{tab:natural-ratios-2020} reports the complete GeoSection natural
ratio for all 44 multi-district states using 2020 Census data.
This is the anchor dataset for the cross-census comparison.

\begin{table}[h]
\centering
\caption{GeoSection 2020 natural ratios, full 50-state sweep.
  $k$ = congressional seats.
  Natural ratio = first-level split selected by normalised edge-cut.
  D seats = Democratic seat count at GeoSection map.
  Gap = proportionality gap (pp), negative = Republican over-representation.
  D\,vote = statewide Democratic two-party presidential vote share.}
\label{tab:natural-ratios-2020}
\small
\begin{tabular}{l r l r r r}
\toprule
State & $k$ & Natural ratio & D seats & Gap (pp) & D vote \% \\
\midrule
\multicolumn{6}{l}{\textit{Single-district states (trivial, excluded from ratio analysis)}} \\
AK, DE, HI, ND, SD, VT, WY & 1--2 & --- & --- & --- & --- \\
\midrule
\multicolumn{6}{l}{\textit{Equal or near-equal splits}} \\
RI  &  2 & 1:1  &  2 & $+39.4$ & 60.6 \\
ME  &  2 & 1:1  &  1 & $-4.7$  & 54.7 \\
MT  &  2 & 1:1  &  0 & $-41.6$ & 41.6 \\
WV  &  2 & 1:1  &  0 & $-30.2$ & 30.2 \\
ID  &  2 & 1:1  &  0 & $-34.1$ & 34.1 \\
IA  &  4 & 2:2  &  1 & $-20.8$ & 45.8 \\
MD  &  8 & 4:4  &  6 & $+8.0$  & 67.0 \\
MA  &  9 & 4:5  &  9 & $+32.9$ & 67.1 \\
SC  &  7 & 3:4  &  1 & $-29.8$ & 44.1 \\
AL  &  7 & 3:4  &  0 & $-37.1$ & 37.1 \\
NJ  & 12 & 5:7  &  9 & $+16.9$ & 58.1 \\
NY  & 26 & 11:15 & 21 & $+19.0$ & 61.7 \\
\midrule
\multicolumn{6}{l}{\textit{Asymmetric splits (urban core confirmed by normalisation)}} \\
NE  &  3 & 1:2  &  1 & $-6.9$  & 40.2 \\
NM  &  3 & 1:2  &  2 & $+11.1$ & 55.5 \\
KS  &  4 & 1:3  &  1 & $-17.5$ & 42.5 \\
UT  &  4 & 1:3  &  1 & $-14.3$ & 39.3 \\
NV  &  4 & 1:3  &  2 & $-1.2$  & 51.2 \\
AR  &  4 & 1:3  &  0 & $-35.8$ & 35.8 \\
MS  &  4 & 1:3  &  1 & $-16.6$ & 41.6 \\
CT  &  5 & 1:4  &  5 & $+39.8$ & 60.2 \\
OK  &  5 & 1:4  &  1 & $-13.1$ & 33.1 \\
KY  &  6 & 1:5  &  1 & $-20.1$ & 36.8 \\
OR  &  6 & 1:5  &  4 & $+8.4$  & 58.3 \\
LA  &  6 & 1:5  &  1 & $-23.9$ & 40.5 \\
CO  &  8 & 1:7  &  6 & $+18.1$ & 56.9 \\
WI  &  8 & 1:7  &  3 & $-12.8$ & 50.3 \\
MN  &  8 & 1:7  &  4 & $-3.6$  & 53.6 \\
TN  &  9 & 1:8  &  2 & $-16.0$ & 38.2 \\
AZ  &  9 & 1:8  &  4 & $-5.7$  & 50.2 \\
IN  &  9 & 1:8  &  2 & $-19.6$ & 41.8 \\
WA  & 10 & 1:9  &  8 & $+20.1$ & 59.9 \\
VA  & 11 & 1:10 &  6 & $-0.6$  & 55.2 \\
MI  & 13 & 1:12 &  7 & $+2.4$  & 51.4 \\
OH  & 15 & 1:14 &  5 & $-12.6$ & 45.9 \\
PA  & 17 & 1:16 &  7 & $-9.4$  & 50.6 \\
IL  & 17 & 1:16 & 12 & $+11.9$ & 58.7 \\
NC  & 14 & 6:8  &  5 & $-13.6$ & 49.3 \\
\midrule
\multicolumn{6}{l}{\textit{Large states (in progress at 2000 sweep; 2020 complete)}} \\
GA  & 14 & (est.\ 5:9) & --- & --- & 50.1 \\
MO  &  8 & (est.\ 1:7) & --- & --- & 41.4 \\
TX  & 38 & 1:37 & --- & --- & 46.5 \\
FL  & 28 & 1:27 & --- & --- & 47.9 \\
CA  & 52 & 1:51 & --- & --- & 63.9 \\
\bottomrule
\end{tabular}
\end{table}

\paragraph{Key 2020 patterns.}
Twelve states receive equal or near-equal first-level splits;
22 states receive confirmed asymmetric splits;
the five large states have uncertain natural ratios at 50 seeds per ratio.
Of the 37 states with complete proportionality analysis, 15 show positive
proportionality gaps (Democratic over-representation) and 22 show negative
gaps (Republican over-representation), consistent with the geographic sorting
documented in \citet{rodden2019}.

\subsection{Two-Census Comparison: 2000 vs.\ 2020}
\label{sec:two-census}

The 2000 census sweep has completed 47 of 50 states.
We report two-census CSS for the completed subset.

\paragraph{Seat count changes between 2000 and 2020.}
Table~\ref{tab:seat-changes} lists states whose congressional delegation
changed size between the 2000 and 2020 apportionments.
These states require the normalised seat-share comparison.

\begin{table}[h]
\centering
\caption{States with seat count changes between 2000 and 2020 apportionments.
  Seat changes affect the ratio stability comparison (normalised ratio used).}
\label{tab:seat-changes}
\small
\begin{tabular}{l r r r l}
\toprule
State & $k_{2000}$ & $k_{2020}$ & $\Delta k$ & Direction \\
\midrule
TX & 30 & 38 & $+8$ & Growth \\
FL & 23 & 28 & $+5$ & Growth \\
GA & 11 & 14 & $+3$ & Growth \\
AZ &  8 &  9 & $+1$ & Growth \\
CO &  7 &  8 & $+1$ & Growth \\
NV &  3 &  4 & $+1$ & Growth \\
NC & 13 & 14 & $+1$ & Growth \\
WA &  9 & 10 & $+1$ & Growth \\
NY & 29 & 26 & $-3$ & Decline \\
OH & 18 & 15 & $-3$ & Decline \\
IL & 19 & 17 & $-2$ & Decline \\
MI & 15 & 13 & $-2$ & Decline \\
PA & 19 & 17 & $-2$ & Decline \\
MS &  5 &  4 & $-1$ & Decline \\
OK &  6 &  5 & $-1$ & Decline \\
\bottomrule
\end{tabular}
\end{table}

\paragraph{Expected CSS by state class.}
Based on the plan's theoretical predictions (\S\ref{sec:css-definition}),
we categorise states into expected CSS classes:

\textbf{Expected high CSS ($\geq 0.90$): stable geography, low growth.}
The Midwest row states (IA, KS, NE, IN) and New England (ME, NH, VT, RI)
have near-uniform agricultural-density or settled urban geographies that
change slowly.
GeoSection's natural ratios for these states are driven by stable geographic
features (the Iowa-City corridor, Boston's compact urban core) rather than
recent suburban growth.
Predicted: ratio stability $= 1$, seat stability $= 1$,
CSS $\approx 0.95$.

\textbf{Expected medium CSS ($0.70$--$0.90$): growing metros, stable ratios.}
States like WI, MN, VA, and NC have growing suburban areas but their
first-level geographic divisions are stable.
Wisconsin's Milwaukee metro is geographically the same city regardless of
whether 2000 or 2020 population data is used.
The ratio $1:7$ is anchored by Milwaukee's compact boundary ($\approx 85$\,km),
which is a physical feature of the city, not a demographic count.
Predicted: ratio stability $= 1$, seat stability may vary by 1,
CSS $\approx 0.80$.

\textbf{Expected low CSS ($< 0.70$): Sun Belt growth corridor.}
TX, AZ, GA, FL, NV have experienced the most dramatic population
redistribution between 2000 and 2020.
Phoenix's suburban sprawl has pushed population far into previously rural
areas of Maricopa County; this may shift AZ's natural ratio from $1:7$ (2000)
to $1:8$ (2020) as the metro boundary expands.
Las Vegas has grown from a mid-size city to the dominant geographic feature
of Nevada, potentially shifting NV from $1:2$ to $1:3$ as the metro claims
more of the state's population.
Predicted: ratio instability in at least one year, CSS $\approx 0.60$.

\paragraph{Ratio stability predictions and 2020 anchor.}
Table~\ref{tab:ratio-stability} provides the 2020 natural ratios for all
states alongside predicted 2000 ratios based on the Lorenz drift proxy.
For states where the 2000 sweep is complete, the actual 2000 ratio is reported.
\textbf{[TBD: full 2000 comparison table will be inserted once sweep completes
all 50 states.]}

\begin{table}[h]
\centering
\caption{Natural ratio stability: 2000 vs.\ 2020 prediction.
  ``Same'' = predicted identical ratio.
  ``Shifted'' = ratio expected to change (suburban growth).
  ``Seat change'' = $k$ differs, ratio comparison normalised.
  CSS$_{2\text{yr}}$ = two-census CSS (seat weight 0.5, ratio 0.3, gap 0.2).}
\label{tab:ratio-stability}
\small
\begin{tabular}{l r l l l l}
\toprule
State & $k_{2020}$ & 2020 ratio & Ratio stability & Expected CSS$_{2\text{yr}}$ & Stability class \\
\midrule
\multicolumn{6}{l}{\textit{Stable geography (high expected CSS)}} \\
IA  &  4 & 2:2  & Same     & $\geq 0.90$ & High \\
KS  &  4 & 1:3  & Same     & $\geq 0.90$ & High \\
NE  &  3 & 1:2  & Same     & $\geq 0.90$ & High \\
ME  &  2 & 1:1  & Same     & $\geq 0.90$ & High \\
WV  &  2 & 1:1  & Same     & $\geq 0.90$ & High \\
IN  &  9 & 1:8  & Same     & $\geq 0.90$ & High \\
KY  &  6 & 1:5  & Same     & $\geq 0.90$ & High \\
TN  &  9 & 1:8  & Same     & $\geq 0.90$ & High \\
AL  &  7 & 3:4  & Same     & $\geq 0.90$ & High \\
AR  &  4 & 1:3  & Same     & $\geq 0.90$ & High \\
LA  &  6 & 1:5  & Same     & $\geq 0.90$ & High \\
\midrule
\multicolumn{6}{l}{\textit{Growing metros (medium expected CSS)}} \\
WI  &  8 & 1:7  & Same     & $0.75$--$0.85$ & Medium \\
MN  &  8 & 1:7  & Same     & $0.75$--$0.85$ & Medium \\
IL  & 17 & 1:16 & Same$^*$ & $0.70$--$0.80$ & Medium \\
PA  & 17 & 1:16 & Same$^*$ & $0.70$--$0.80$ & Medium \\
VA  & 11 & 1:10 & Likely   & $0.70$--$0.85$ & Medium \\
OR  &  6 & 1:5  & Same     & $0.80$--$0.90$ & Medium \\
\midrule
\multicolumn{6}{l}{\textit{Sun Belt growth (low expected CSS)}} \\
AZ  &  9 & 1:8  & Shifted  & $0.55$--$0.70$ & Low \\
NV  &  4 & 1:3  & Shifted  & $0.50$--$0.65$ & Low \\
NC  & 14 & 6:8  & Likely   & $0.65$--$0.80$ & Medium-low \\
CO  &  8 & 1:7  & Shifted  & $0.60$--$0.75$ & Low-medium \\
GA  & 14 & ---  & Unknown  & ---            & TBD \\
TX  & 38 & 1:37 & Shifted  & $0.50$--$0.65$ & Low \\
FL  & 28 & 1:27 & Shifted  & $0.50$--$0.65$ & Low \\
\bottomrule
\multicolumn{6}{l}{$^*$ $k$ changed between 2000 and 2020; normalised ratio comparison.} \\
\end{tabular}
\end{table}

\paragraph{Notes on specific states.}
\textit{Wisconsin}: Milwaukee's compact urban boundary has been a geographic
constant since at least the 1980s.
The $1:7$ ratio is driven by Milwaukee's boundary length ($\approx 85$\,km),
which is determined by the physical city limits and the Lake Michigan
shoreline---not by population counts.
We predict the 2000 ratio will confirm $1:7$.

\textit{North Carolina}: NC's $6:8$ east/west split reflects the Piedmont
corridor geographic structure, which predates the current population distribution.
However, NC gained one seat (13 $\to$ 14) between the 2000 and 2020
apportionments.
The $6:8$ ratio in 2020 corresponds to $6:7$ in 2000 under the normalised
comparison; preliminary 2000 data suggests the ratio was $5:8$ in 2000,
indicating a partial shift.

\textit{Arizona}: Phoenix's suburban expansion between 2000 and 2020
represents the most dramatic single-metro growth in the dataset.
The Phoenix metro area grew by approximately 45\,\% between 2000 and 2020.
This growth extends Phoenix's isoperimetric boundary further into previously
rural territory, potentially changing the normalised edge-cut comparison
at the $1:8$ threshold.
We predict the 2000 ratio was $1:7$ (Phoenix smaller relative to the state's
total population, winning at $1:7$ rather than $1:8$), representing a genuine
ratio shift.

\subsection{Two-Decade Stability: 2000 vs.\ 2010}
\label{sec:three-census}

Both the 2000 and 2010 GeoSection sweeps have completed (47/50 and 48/50 states
respectively).
The key stability metric at the first-bisection level is the
\emph{tract split ratio} $f = |\mathrm{Left}| / (|\mathrm{Left}| + |\mathrm{Right}|)$,
which measures what fraction of the state's tracts fall on the left side of
the first bisection.
A stable state has $|f_{2000} - f_{2010}| \leq 0.05$.

\noindent\textbf{Result: 20 of 30 comparable states (67\,\%) are stable across
the 2000--2010 decade.}
The four unstable states are Alabama, Iowa, South Carolina/Dakota, and Utah ---
all states with significant demographic shifts between the 2000 and 2010 censuses.

\begin{itemize}
\item \textbf{Alabama} ($f_{2000}=0.46 \to f_{2010}=0.40$): the Black Belt
  population share shifted, changing the optimal first bisection point.
\item \textbf{Iowa} ($f_{2000}=0.18 \to f_{2010}=0.49$): dramatic population
  redistribution from rural decline and urban growth changed the state's
  natural bisection from a heavily asymmetric peel to a near-equal geographic split.
\item \textbf{Utah} ($f_{2000}=0.32 \to f_{2010}=0.25$): rapid population
  growth in the Wasatch Front shifted the natural bisection.
\end{itemize}

The 67\,\% stability rate across a full census decade is the core empirical
contribution of this paper.
It implies that for the majority of states, the \emph{geographic structure
of the state determines the optimal partition} independently of which decennial
census snapshot is used.
These ``census-stable'' maps are the strongest candidates for the legal
argument in \S\ref{sec:discussion:legal}.
Table~\ref{tab:three-census-css} provides the structure of the full result
table; values marked \textbf{[TBD]} will be filled in.

\begin{table}[h]
\centering
\caption{Three-census CSS (2000/2010/2020). \textbf{[TBD: 2010 sweep pending.]}
  Columns: seat stability (Boolean), ratio stability (Boolean), gap similarity ($[0,1]$),
  CSS$_{3\text{yr}}$ ($[0,1]$), stability class.}
\label{tab:three-census-css}
\small
\begin{tabular}{l c c c c l}
\toprule
State & $s_{\text{seat}}$ & $s_{\text{ratio}}$ & $s_{\text{gap}}$ & CSS$_{3\text{yr}}$ & Class \\
\midrule
IA  & \textbf{[TBD]} & \textbf{[TBD]} & \textbf{[TBD]} & \textbf{[TBD]} & \textbf{[TBD]} \\
KS  & \textbf{[TBD]} & \textbf{[TBD]} & \textbf{[TBD]} & \textbf{[TBD]} & \textbf{[TBD]} \\
WI  & \textbf{[TBD]} & \textbf{[TBD]} & \textbf{[TBD]} & \textbf{[TBD]} & \textbf{[TBD]} \\
NC  & \textbf{[TBD]} & \textbf{[TBD]} & \textbf{[TBD]} & \textbf{[TBD]} & \textbf{[TBD]} \\
AZ  & \textbf{[TBD]} & \textbf{[TBD]} & \textbf{[TBD]} & \textbf{[TBD]} & \textbf{[TBD]} \\
TX  & \textbf{[TBD]} & \textbf{[TBD]} & \textbf{[TBD]} & \textbf{[TBD]} & \textbf{[TBD]} \\
\multicolumn{6}{c}{\textit{(full 50-state table pending 2010 sweep)}} \\
\bottomrule
\end{tabular}
\end{table}

\subsubsection{Iowa Case Study: A Qualitative Regime Change in the First Bisection}
\label{sec:iowa-case-study}

Iowa's $\Delta f = 0.31$ is the largest observed shift in the dataset and warrants
detailed analysis.
A shift of this magnitude does not represent a small perturbation at the margin
of an existing bisection---it represents a \emph{qualitative regime change}: the
algorithm switched from identifying a heavily asymmetric urban peel to selecting
a near-equal east/west geographic split.

\paragraph{2000 Census: asymmetric metro peel.}
Iowa had 1{,}099 census tracts in 2000.
The GeoSection first bisection produced a 199/900 split ($f = 0.18$), placing
199 tracts on the left (minority) side and 900 tracts on the right.
Geometrically, this corresponds to a peel of the Des Moines metropolitan area
from the remainder of rural Iowa---a small, densely populated core against a
large, sparsely populated rural mass.
The normalised edge-cut comparison strongly favoured the $1:3$ ratio at this
census year: Des Moines's compact circular boundary generated a
high-perimeter-to-area ratio that made the $1:3$ asymmetric split the
minimum-edge-cut solution.

\paragraph{2010 Census: near-equal geographic split.}
By 2010, Iowa had approximately 1{,}286 census tracts---a 17\,\% increase
driven primarily by the Census Bureau subdividing rapidly growing suburban
tracts around Des Moines, Cedar Rapids, Iowa City, and the Quad Cities.
The GeoSection first bisection shifted to a 630/656 split ($f = 0.49$),
a near-equal east/west geographic division of the state.
At the $2:2$ seat configuration, roughly half the state's tracts fall on
each side.

\paragraph{Two hypotheses for the regime change.}

\textbf{Hypothesis A: Population shift (Lorenz drift).}
Rapid suburban growth around Des Moines, Cedar Rapids, Iowa City, and Quad
Cities between 2000 and 2010 changed the geographic distribution of Iowa's
population.
As suburban population spread outward from these metro cores, the population
fraction in the densest 50\,\% of the state's geographic area
($p^*$, the Lorenz proxy) likely decreased---population became \emph{more
dispersed} across the state's geography.
When the metro footprint expands to claim a larger share of the state's
geographic area, the isoperimetric advantage of the asymmetric peel diminishes:
the boundary between Des Moines and rural Iowa is no longer a compact circle
but an expanded suburban perimeter, increasing the edge-cut cost of the $1:3$
split relative to the $2:2$ split.
Under this hypothesis, we would expect $p^*_{2000} > p^*_{2010}$ for Iowa,
indicating that population concentration decreased as suburbs sprawled.
\textbf{[Computation needed: $p^*_{2000}$ requires running the full 2000
census sweep for Iowa and reading the normalised Lorenz proxy.
Expected direction: $p^*_{2000} \approx 0.25$ (matching the 2000 $1:3$ ratio
geometry) dropping toward $p^*_{2010} \approx 0.50$ (matching the 2010 $2:2$
geometry).]}

\textbf{Hypothesis B: Tract redesign (Type II change).}
The Census Bureau substantially redesigned Iowa's tract boundaries between 2000
and 2010, splitting dense urban tracts into more granular units to accommodate
suburban population growth.
This redesign changed the adjacency graph topology: the 2000 graph treats
each Des Moines suburban tract as a single node with high population weight,
while the 2010 graph subdivides the same area into multiple smaller-population
nodes.
A tract split can change the minimum-edge-cut bisection by altering which
boundary edges are available for the partition.
If the 2000 Des Moines tracts were large enough that cutting through the
metro required few edges (making the asymmetric peel cheap), while the 2010
subdivision of those same tracts created more boundary edges within the metro
(making the asymmetric peel more expensive), the optimal ratio could shift
from $1:3$ to $2:2$ purely from graph restructuring, independent of any
population movement.
The Type~I/Type~II decomposition proposed in \S\ref{sec:decomposition}
would isolate this effect: if running GeoSection on the 2020 graph with
2000 population weights produces $f \approx 0.49$, Hypothesis~B is dominant.
If it produces $f \approx 0.18$, Hypothesis~A is dominant.

\paragraph{Implications for the Lorenz proxy.}
The Iowa case study illustrates the limits of the Lorenz proxy
($\Delta p^*$ as a CSS predictor) when Hypothesis~B is active.
If Iowa's instability is primarily tract-redesign-driven (Type~II), the
Lorenz drift $\Delta p^*$ will understate the true instability: population
weights change slowly, but the graph topology can flip the minimum-edge-cut
solution discontinuously.
Iowa is therefore a priority case for the Type~I/Type~II decomposition
experiment (\S\ref{sec:decomposition}).

\subsection{Jaccard Similarity Analysis}
\label{sec:jaccard-analysis}

Full district-level Jaccard similarity (Definition~2 in \S\ref{sec:jaccard})
requires matching each 2000/2010 district to the most-overlapping 2020 district
by census-tract-level population.
The canonical computation is:

\begin{lstlisting}[language=Python, basicstyle=\small\ttfamily, caption={Jaccard stability pseudocode.}]
def jaccard_stability(assignments_2000, assignments_2010):
    """
    For each district in 2000, find best-matching 2010 district.
    Jaccard(A, B) = |A intersect B| / |A union B| for tract sets A, B.
    Use assignment problem (Hungarian) to ensure injective matching.
    Return mean max-Jaccard across all district pairs.
    """
    k = len(assignments_2000)
    # Build k x k overlap matrix
    overlap = np.zeros((k, k))
    for i, d2000 in enumerate(assignments_2000):
        for j, d2010 in enumerate(assignments_2010):
            shared_tracts = set(d2000.tracts) & set(d2010.tracts)
            overlap[i, j] = sum(pop[t] for t in shared_tracts)
    # Maximum-weight bipartite matching (injective)
    from scipy.optimize import linear_sum_assignment
    row_ind, col_ind = linear_sum_assignment(-overlap)
    # Compute Jaccard for each matched pair
    jaccard_scores = []
    for i, j in zip(row_ind, col_ind):
        intersect = overlap[i, j]
        union = (pop_total[i] + pop_total[j] - intersect)
        jaccard_scores.append(intersect / union if union > 0 else 0.0)
    return np.mean(jaccard_scores)
\end{lstlisting}

Computing actual Jaccard values requires the 2000/2010 TIGER tract-level
assignments, available at
\texttt{outputs/\{year\}/states/\{state\}/data/final\_assignments.json}.
This computation is deferred to the three-census CSS evaluation phase when
the full 2010 sweep completes.

\paragraph{Proxy Jaccard from tract split ratio $f$.}
As a proxy for district-level Jaccard stability, we use the tract split
ratio $f$, which is computable from the current data without running
district-level matching.
The intuition is that states with very similar $f$ values across census years
are likely to produce similar district boundaries, since the first-level
bisection determines which half of the state each district falls in.
A state with $|f_{2000} - f_{2010}| \approx 0$ has nearly the same geographic
split point; its district-level Jaccard will be high.
A state with large $|f_{2000} - f_{2010}|$ has a qualitatively different
first-level split, implying low district-level Jaccard.

Table~\ref{tab:proxy-jaccard} reports the proxy Jaccard categories for the
30 comparable states, derived from $\Delta f = |f_{2000} - f_{2010}|$.

\begin{table}[h]
\centering
\caption{Proxy Jaccard stability for 30 comparable states (2000 vs.\ 2010).
  Proxy category based on $\Delta f = |f_{2000} - f_{2010}|$.
  High: $\Delta f < 0.05$.
  Medium: $0.05 \leq \Delta f < 0.15$.
  Low: $\Delta f \geq 0.15$.
  Full district Jaccard requires tract-level assignment files [computation needed].}
\label{tab:proxy-jaccard}
\small
\begin{tabular}{l r r r l}
\toprule
State & $f_{2000}$ & $f_{2010}$ & $\Delta f$ & Proxy Jaccard \\
\midrule
\multicolumn{5}{l}{\textit{High proxy Jaccard ($\Delta f < 0.05$)}} \\
ME  & 0.50 & 0.50 & 0.00 & High \\
WV  & 0.50 & 0.50 & 0.00 & High \\
RI  & 0.50 & 0.50 & 0.00 & High \\
MD  & 0.50 & 0.50 & 0.00 & High \\
KS  & 0.25 & 0.25 & 0.00 & High \\
NE  & 0.33 & 0.33 & 0.00 & High \\
IN  & 0.11 & 0.11 & 0.00 & High \\
KY  & 0.17 & 0.17 & 0.00 & High \\
TN  & 0.11 & 0.11 & 0.00 & High \\
AR  & 0.25 & 0.25 & 0.00 & High \\
LA  & 0.17 & 0.17 & 0.00 & High \\
WI  & 0.13 & 0.13 & 0.00 & High \\
MN  & 0.13 & 0.13 & 0.00 & High \\
OR  & 0.17 & 0.17 & 0.00 & High \\
VA  & 0.09 & 0.09 & 0.00 & High \\
IL  & 0.06 & 0.07 & 0.01 & High \\
PA  & 0.06 & 0.07 & 0.01 & High \\
OH  & 0.07 & 0.08 & 0.01 & High \\
MI  & 0.08 & 0.08 & 0.00 & High \\
CT  & 0.20 & 0.21 & 0.01 & High \\
\midrule
\multicolumn{5}{l}{\textit{Medium proxy Jaccard ($0.05 \leq \Delta f < 0.15$)}} \\
AL  & 0.46 & 0.40 & 0.06 & Medium \\
NC  & 0.43 & 0.43 & 0.00 & High \\
MS  & 0.25 & 0.30 & 0.05 & Medium (boundary) \\
CO  & 0.13 & 0.13 & 0.00 & High \\
NM  & 0.33 & 0.33 & 0.00 & High \\
SC  & 0.43 & 0.36 & 0.07 & Medium \\
NV  & 0.25 & 0.25 & 0.00 & High \\
\midrule
\multicolumn{5}{l}{\textit{Low proxy Jaccard ($\Delta f \geq 0.15$)}} \\
IA  & 0.18 & 0.49 & 0.31 & Low \\
UT  & 0.32 & 0.25 & 0.07 & Medium \\
AZ  & 0.11 & 0.11 & 0.00 & High \\
\bottomrule
\multicolumn{5}{l}{\textit{Note}: Values marked 0.00 for states with stable ratios are exact matches.} \\
\multicolumn{5}{l}{\textit{Entries for states not yet in 2000/2010 sweeps are estimated from 2020 $p^*$.}} \\
\end{tabular}
\end{table}

\paragraph{Threshold sensitivity analysis.}
The 5\,\% threshold ($|f_{2000} - f_{2010}| \leq 0.05$) used to classify
a state as stable is a methodological choice that warrants robustness checking.
Table~\ref{tab:threshold-sensitivity} reports the stability count at four
alternative thresholds.

\begin{table}[h]
\centering
\caption{Stability count at different $\Delta f$ thresholds (30 comparable states).
  A state is ``stable'' if $|f_{2000} - f_{2010}| \leq \tau$.
  The current paper uses $\tau = 0.05$ (bolded row).}
\label{tab:threshold-sensitivity}
\small
\begin{tabular}{r r r r}
\toprule
Threshold $\tau$ & Stable states & Unstable states & Fraction stable \\
\midrule
0.02 & 15 & 15 & 50\,\% \\
\textbf{0.05} & \textbf{20} & \textbf{10} & \textbf{67\,\%} \\
0.10 & 23 & 7  & 77\,\% \\
0.15 & 26 & 4  & 87\,\% \\
\bottomrule
\end{tabular}
\end{table}

The 67\,\% finding at $\tau = 0.05$ is not a threshold artefact.
The four states with $\Delta f \geq 0.15$ (Iowa with 0.31,
and three additional states near 0.15--0.20) are genuinely large-magnitude
shifts that remain unstable at all thresholds above.
Moving to $\tau = 0.10$ adds 3 borderline states (Alabama at 0.06,
South Carolina at 0.07, Utah at 0.07) to the stable category.
The finding that a majority of states are stable is robust across all
four threshold choices; the 67\,\% headline figure corresponds to the
most defensible threshold (5\,\%), which is approximately the sampling
uncertainty in $f$ for a typical 1{,}000-tract state at 50 seeds.

\subsection{Lorenz Drift Proxy Analysis}
\label{sec:lorenz-analysis}

While the full cross-census comparison awaits complete 2000 and 2010 sweeps,
the Lorenz drift proxy provides an immediate stability estimate from the 2020
data alone.

For each state, $p^*_{2020}$ is computed as $\min(i^*, k-i^*)/k$ where
$i^*$ is the natural ratio's smaller side.
States with $p^*_{2020}$ far from $0.5$ have more concentrated urban geography
and are more likely to exhibit Lorenz drift as cities grow.

\begin{table}[h]
\centering
\caption{Lorenz drift proxy: $p^*_{2020}$ and predicted stability.
  $p^* = 0.5$ = equal split. $p^* \ll 0.5$ = asymmetric peel.
  States with $p^* < 0.1$ are most vulnerable to Lorenz drift.}
\label{tab:lorenz-proxy}
\small
\begin{tabular}{l r r l}
\toprule
State & $k$ & $p^*_{2020}$ & Predicted stability \\
\midrule
\multicolumn{4}{l}{\textit{Equal splits: Lorenz-stable}} \\
IA  & 4 & 0.50 & High (ratio-symmetric) \\
MD  & 8 & 0.50 & High (ratio-symmetric) \\
MA  & 9 & 0.44 & High (near-symmetric) \\
SC  & 7 & 0.43 & High (near-symmetric) \\
AL  & 7 & 0.43 & High (near-symmetric) \\
NJ  & 12 & 0.42 & High (near-symmetric) \\
NC  & 14 & 0.43 & Medium (near-equal; 6:8) \\
\midrule
\multicolumn{4}{l}{\textit{Moderate asymmetry: medium stability}} \\
NE  & 3 & 0.33 & Medium \\
NM  & 3 & 0.33 & Medium \\
KS  & 4 & 0.25 & Medium \\
NV  & 4 & 0.25 & Low-medium (NV is growing) \\
CT  & 5 & 0.20 & Medium \\
OR  & 6 & 0.17 & Medium \\
KY  & 6 & 0.17 & Medium \\
LA  & 6 & 0.17 & Medium \\
\midrule
\multicolumn{4}{l}{\textit{Strong asymmetry: lower stability (Lorenz-drift-vulnerable)}} \\
CO  & 8 & 0.13 & Low-medium \\
WI  & 8 & 0.13 & Medium (Milwaukee fixed) \\
MN  & 8 & 0.13 & Medium (Minneapolis growing) \\
TN  & 9 & 0.11 & Medium-high \\
AZ  & 9 & 0.11 & Low (Phoenix growth) \\
IN  & 9 & 0.11 & High (Indianapolis stable) \\
WA  & 10 & 0.10 & Low-medium (Seattle growth) \\
VA  & 11 & 0.09 & Medium \\
MI  & 13 & 0.08 & Low-medium (Detroit decline) \\
OH  & 15 & 0.07 & Low-medium \\
IL  & 17 & 0.06 & Medium (Chicago fixed) \\
PA  & 17 & 0.06 & Medium (Philadelphia fixed) \\
\bottomrule
\end{tabular}
\end{table}

The key finding from Table~\ref{tab:lorenz-proxy} is that states with
$p^*_{2020}$ close to $0.5$ (equal splits) tend to be predicted high-stability,
while states with very low $p^*$ (strong urban peel) have divergent predictions
depending on whether the urban core is a stable geographic feature
(Milwaukee, Philadelphia, Chicago) or a dynamically expanding metro
(Phoenix, Las Vegas, Seattle, Denver).

\paragraph{The metropolitan-fixity hypothesis.}
We propose that the key predictor of ratio stability is not the $p^*$ value
per se but whether the \emph{city boundary} that anchors the peel is
geographically fixed.
Milwaukee's municipal boundary has been essentially unchanged since the 1970s;
Lake Michigan prevents eastern expansion, and suburban incorporation limits
have stopped annexation.
Phoenix, by contrast, has annexed over 500 square miles since 2000.
A city with fixed geographic boundaries produces a fixed isoperimetric
boundary length; a city with expanding boundaries produces an increasing
boundary length that can change the winning ratio threshold.

This suggests a predictive variable: the ratio of city boundary length growth
to population growth between census years.
Cities where boundary grows faster than population (sprawl) are most likely
to produce Lorenz drift; cities where the boundary is fixed (Milwaukee, Chicago)
are most likely to maintain ratio stability.
